\documentclass{article}
\usepackage{amsmath}
\usepackage{amssymb}
\usepackage{fullpage}
\usepackage{enumerate}
\usepackage{xcolor}
\usepackage{bm}
\pagestyle{empty}
\usepackage{graphicx}
\graphicspath{{C:\Users\steve\Pictures}}
\begin{document}

\noindent{{\Large \bf Fall 25, 4100 Problem Set 7, Hayden Fuller}

\large

\bigskip
\noindent{Please complete all problems. For any proof questions, please structure your proof similar to those from lecture.

\begin{enumerate}

\item Let $A\in\mathbb{C}^{n\times n}$ and $W$ be a subspace of $\mathbb{C}^n$. Recall that $W$ is an $A$-invariant subspace if and only if for all ${\bm x}\in W$, $A{\bm x}\in W$. Prove that the generalized eigenspace, ${\rm null}(A-\lambda I)^n$ is $A$-invariant for any eigenvalue $\lambda$ and $n$.
\\
\\let $x\in\textrm{null}(A-\lambda I)^n$
\\by definition $(A-\lambda I)^nx=0$
\\$A=(A-\lambda I)+\lambda I$
\\$Ax=((A-\lambda I)+\lambda I)x$
\\because both are polynomials in A they commute $(A-\lambda I)A=A(A-\lambda I)$
\\$(A-\lambda I)^nA=A(A-\lambda I)^n$
\\$(A-\lambda I)^n(Ax)=A(A-\lambda I)^nx$
\\but $(A-\lambda I)^nx=0$
\\so $A(A-\lambda I)^nx=A\cdot0=0$
\\and $(A-\lambda I)^n(Ax)=0$
\\therefore $Ax\in\textrm{null}(A-\lambda I)^n$
\\QED

\newpage\item For each of the following matrices find the eigenvalues, eigenvectors, and generalized eigenvectors, and state the Jordan form
\begin{enumerate}[a.]
\item $\left[\begin{array}{cc}
-3 & -1\\
1 & -5
\end{array}\right]$
\\
\\$\textrm{det}(A-\lambda I)=\textrm{det}\left[\begin{array}{cc} -3-\lambda & -1 \\ 1 & -5-\lambda \end{array}\right]=(-3-\lambda)(-5-\lambda)+1=\lambda^2+8\lambda+16=(\lambda+4)^2$
\\$\lambda=-4$ with algebraic multiplicity 2
\\$(A+4I)x=0$
\\$A+4I=\left[\begin{array}{cc} 1 & -1 \\ 1 & -1 \end{array}\right]$
\\one dimensional null space $x_1=x_2$
\\eigenvector $v_1=\left[\begin{array}{c} 1 \\ 1 \end{array}\right]$
\\geometric multiplicity = 1
\\$(A+4I)v_2=\left[\begin{array}{c} 1 \\ 1 \end{array}\right]$
\\$v_2=\left[\begin{array}{c} 1 \\ 0 \end{array}\right]$
\\$J=\left[\begin{array}{cc} -4 & 1 \\ 0 & -4 \end{array}\right]$
\\
\item $\left[\begin{array}{ccc}
2 & 1 & 3\\
-1 & 0 & -5\\
0 & 0 & 3
\end{array}\right]$
\\
\\$\textrm{det}(B-\lambda I)=\textrm{det}\left[\begin{array}{ccc} 2-\lambda & 1 & 3 \\ -1 & 0-\lambda & -5 \\ 0 & 0 & 3-\lambda \end{array}\right]=(2-\lambda)(-\lambda)(3-\lambda)-0+0-(1)(-1)(3-\lambda)+0-0=(2-\lambda)(-\lambda)(3-\lambda)+(3-\lambda)=-\lambda^3+5\lambda^2-7\lambda+3=(\lambda-1)^2(\lambda-3)$
\\$\lambda=1,3$ with algebraic multiplicities 2 and 1 respectively
\\$(B-I)x=0$
\\$B-I=\left[\begin{array}{ccc} 1 & 1 & 3 \\ -1 & -1 & -5 \\ 0 & 0 & 2 \end{array}\right]$
\\$(B-I)u_1=0$
\\$u_1=\left[\begin{array}{c} -1 \\ 1 \\ 0 \end{array}\right]$
\\$(B-3I)x=0$
\\$B-3I=\left[\begin{array}{ccc} -2 & 1 & 3 \\ -1 & -3 & -5 \\ 0 & 0 & 0 \end{array}\right]$
\\$(B-3I)u_2=0$
\\$u_2=\left[\begin{array}{c} 1 \\ -2 \\ 1 \end{array}\right]$
\\$(B-I)w=u_1$
\\$w=\left[\begin{array}{c} -1 \\ 0 \\ 0 \end{array}\right]$
\\$J=\textrm{diag}(\left[\begin{array}{cc} 1 & 1 \\ 0 & 1 \end{array}\right],\left[\begin{array}{c} 3 \end{array}\right])$
\\$J=\left[\begin{array}{ccc} 1 & 1 & 0 \\ 0 & 1 & 0 \\ 0 & 0 & 3 \end{array}\right]$
\end{enumerate}



\newpage\item \begin{enumerate}[a.]
\item Given $J=\left[\begin{array}{cc}
\lambda & 1\\
0 & \lambda
\end{array}\right]$, find a formula for the matrix exponential $\displaystyle{e^{Jt}}$.
\\
\\let $N=\left[\begin{array}{cc} 0 & 1 \\ 0 & 0 \end{array}\right]$ and write $J=\lambda I + N$
\\note $N^2=0$ and $\lambda I$ commutes with $N$
\\$e^{Jt}=e^{(\lambda I+N)t}=e^{\lambda  t}e^{Nt}$
\\$e^{Nt}=I+Nt$
\\$e^{Jt}=e^{\lambda t}\left[\begin{array}{cc} 1 & t \\ 0 & 1 \end{array}\right]$
\\
\item Use the result above and 2a. to solve the linear system 
$$\begin{array}{rcl}
x_1' &=& -3x_1-x_2 \\
x_2' &=& x_1-5x_2
\end{array}$$
with inital values $x_1(0)=1$, $x_2(0)=2$
\\
\\let $A=\left[\begin{array}{cc} -3 & -1 \\ 1 & -5 \end{array}\right]$ and write $x'=Ax$
\\A is similar to $J=\left[\begin{array}{cc} -4 & 1 \\ 0 & -4 \end{array}\right]$ with similarity matrix $P=\left[\begin{array}{cc} 1 & 1 \\ 1 & 0 \end{array}\right]$, $P^{-1}=\left[\begin{array}{cc} 0 & 1 \\ 1 & -1 \end{array}\right]$
\\$A=PJP^{-1}$ so $e^{At}=Pe^{Jt}P^{-1}$
\\$e^{Jt}=e^{-4t}\left[\begin{array}{cc} 1 & t \\ 0 & 1 \end{array}\right]$
\\$e^{At}=P(e^{-4t}\left[\begin{array}{cc} 1 & t \\ 0 & 1 \end{array}\right])P^{-1}=e^{-4t}P\left[\begin{array}{cc} 1 & t \\ 0 & 1 \end{array}\right]P^{-1}$
\\$e^{At}=e^{-4t}\left[\begin{array}{cc} 1+t & -t \\ t & 1-t \end{array}\right]$
\\$x(t)=e^{At}x(0)=e^{-4t}\left[\begin{array}{cc} 1+t & -t \\ t & 1-t \end{array}\right]\left[\begin{array}{c} 1 \\ 2 \end{array}\right]=e^{-4t}\left[\begin{array}{c} 1+t-2t \\ t+2(1-t) \end{array}\right]=e^{-4t}\left[\begin{array}{c} 1-t \\ 2-t \end{array}\right]$
\\$x_1(t)=(1-t)e^{-4t}$
\\$x_2(t)=(2-t)e^{-4t}$
\end{enumerate}


\newpage\item Let $A$ be such that $p_A(\lambda)=|A-\lambda I|=(\lambda-4)(\lambda-2)^5(\lambda+1)^2$ with $\textrm{nullity}(A-4I)=1$, $\textrm{nullity}(A-2I)=3$, and $\textrm{nullity}(A+I)=1$. Fnd all possible Jordan forms of $A$ up to reordering.
\\
\\$\lambda=4$ a mult $=1$
\\$\lambda=2$ a mult $=5$
\\$\lambda=-1$ a mult $=2$
\\$\textrm{nullity}(A-4I)=1$ one Jordan block for $\lambda=4$
\\$J_1(4)$
\\$\textrm{nullity}(A-2I)=3$ one Jordan block for $\lambda=2$, can be $3+1+1$ or $2+2+1$
\\$J_3(2)\oplus J_1(2) \oplus J_1(2)$ or $J_2(2)\oplus J_2(2) \oplus J_1(2)$
\\$\textrm{nullity}(A+I)=1$ one Jordan block for $\lambda=-1$
\\$J_2(-1)$
\\two possible Jordan forms (up to reordering of blocks)
\\$J=J_1(4) \oplus (J_3(2)\oplus J_1(2) \oplus J_1(2)) \oplus J_2(-1)$
\\$J=J_1(4) \oplus (J_2(2)\oplus J_2(2) \oplus J_1(2)) \oplus J_2(-1)$

\newpage\item Let $A$ be such that $p_A(\lambda)=|A-\lambda I|=(\lambda-3)^7(\lambda+2)^6$ with minimal polynomial, $\mu_A(x)=(x-3)^5(x+2)^3$ and geometric multiplicities, $\textrm{nullity}(A-3I)=2$ and $\textrm{nullity}(A+2I)=3$. Find the {\bf unique} (up to reordering) Jordan form of $A$.
\\
\\$\lambda=7$, a mult $7$, g mult $2$, largest block size $5$, must be $5+2$
\\$J_5(3)$ and $J_2(3)$
\\$\lambda=-2$, a mult $6$, g mult $3$, largest block size $3$, must be $3+2+1$
\\$J_3(-2)$, $J_2(-2)$, and $J_1(-2)$
\\$J=\textrm{diag}(J_5(3), J_2(3), J_3(-2), J_2(-2), J_1(-2))$

\newpage\item Let $A\in\mathbb{R}^{n\times n}$ be such that $A^k=I$. Prove that $A$ is diagonalizable. (Hint: a polynomial $p$ has a simple zero at $z_0$ iff $p(z_0)=0$ but $p'(z_0)\neq 0$.)
\\
\\set $p(x)=x^k-1$
\\then $p(A)=A^k-I=0$
\\so the minimal polynomail $\mu_A(x)$ of A divides $p(x)$
\\the roots of $p(x)$ are the kth roots of unity $\zeta$
\\for each root we have $p'(\zeta)=k\zeta^{k-1}$
\\$\zeta\neq0$ and $k\neq0$ so $p'(\zeta)\neq0$
\\by the hint, every root of $p$ is a simple root, no repeated factors
\\therefore square free, so any divisor including $\mu_A(x)$ is also square free
\\therefore the minimal polynomial has one simple distinct linear factors
\\therefore $A$ is diagonalizable
\\QED













\end{enumerate}


\end{document}